\documentclass[a4paper, 12pt]{article}
\usepackage[hidelinks]{hyperref} % hidelinks remove as bordas coloridas dos links
% --- Codificação e Idioma ---
\usepackage[T1]{fontenc}    % Essencial para acentos e caracteres corretos
\usepackage[utf8]{inputenc}
\usepackage[portuguese]{babel}

% --- Fontes e Estilo ---
\usepackage{mathptmx}      % Times New Roman
\usepackage{enumitem}
\usepackage{indentfirst}   % Indenta o primeiro parágrafo de cada seção
\usepackage{geometry}
\geometry{left=30mm,right=20mm,top=30mm,bottom=20mm,includefoot}

% --- Imagens ---
\usepackage{graphicx}
\usepackage{float}         % Ajuda a fixar imagens com [H]
\graphicspath{ {./images/} }

\usepackage{tikz}
\usetikzlibrary{shapes.geometric, arrows.meta, positioning}

% Definição dos estilos dos blocos do fluxograma
\tikzstyle{startstop} = [rectangle, rounded corners, minimum width=3cm, minimum height=1cm, text centered, draw=black, fill=red!20]
\tikzstyle{io} = [trapezium, trapezium left angle=70, trapezium right angle=110, minimum width=3cm, minimum height=1cm, text centered, draw=black, fill=blue!20]
\tikzstyle{process} = [rectangle, minimum width=3.5cm, minimum height=1cm, text centered, draw=black, fill=orange!20]
\tikzstyle{decision} = [diamond, aspect=2, minimum width=3cm, minimum height=1cm, text centered, draw=black, fill=green!20]
\tikzstyle{arrow} = [thick,->,>=stealth]

% --- Títulos (Configuração correta para Sumário e Numeração) ---
\usepackage{titlesec}
\titleformat{\section}{\bfseries\Large}{\thesection}{1em}{}
\titlespacing*{\section}{0pt}{20pt}{10pt}

% --- Cabeçalho e Rodapé ---
\usepackage{fancyhdr}
\pagestyle{fancy}
\fancyhf{}
\renewcommand{\headrulewidth}{0pt}
\rfoot{\thepage}

% --- Espaçamento e Parágrafo ---
\linespread{1.50}
\setlength\parindent{10mm}

% --- Comandos Personalizados ---
\newcommand{\jump}{\vspace{20pt}}

% Ambiente para Citações Longas
\newenvironment{leftquote}[1][]
{
    \hspace{30pt} #1
    \vspace{20pt}
    \fontsize{10}{10}\selectfont
    \begin{flushright}
    \begin{minipage}{120mm}
}
{
    \end{minipage}
    \end{flushright}
    \par
    \vspace{20pt}
}

% Ambiente para Imagens/Tabelas com Fonte
\newcommand{\tmp}{}
\newenvironment{container}[3]
{
    \renewcommand{\tmp}{#3} 
    \begin{center}
    \begin{minipage}{#1}
    \centering
    #2
}
{
    \raggedright
    \vspace{2pt}
    {\footnotesize \tmp}
    \end{minipage}
    \end{center}
}

% ---------------------------------------------------